% !TEX root = ../main.tex
%% Données utiles (constantes, couples, énergies de liaison, tableau périodique) et repères du programme de première

\editionPage{donnees-utiles}{Données utiles}{Constantes, conversions et données de référence}{Valeurs arrondies, suffisantes pour les exercices · Constantes fondamentales d'après le CODATA}
\begin{blocEdition}{donnees-utiles}
\fontsize{10}{13.5}\selectfont
\editionEtape{Constantes physiques}
{\small\begin{tabularx}{\linewidth}{@{}Yl|Yl@{}}\toprule
Célérité de la lumière dans le vide & $c=\qty{3,00e8}{m.s^{-1}}$ & Constante d'Avogadro & $N_\mathrm{A}=\qty{6,02e23}{mol^{-1}}$ \\
Constante de Planck & $h=\qty{6,63e-34}{J.s}$ & Constante de gravitation & $G=\qty{6,67e-11}{N.m^2.kg^{-2}}$ \\
Charge élémentaire & $e=\qty{1,60e-19}{C}$ & Constante de la loi de Coulomb & $k=\qty{9,0e9}{N.m^2.C^{-2}}$ \\
Masse de l'électron & $m_e=\qty{9,11e-31}{kg}$ & Masses du proton et du neutron & $m_p\approx m_n=\qty{1,67e-27}{kg}$ \\
Pesanteur (France) & $g=\qty{9,81}{N.kg^{-1}}$ & Pression atmosphérique normale & $\qty{1,013e5}{Pa}$ \\\bottomrule
\end{tabularx}}
\editionEtape{Conversions et données courantes}
{\small\begin{tabularx}{\linewidth}{@{}Yl|Yl@{}}\toprule
Électronvolt & $\qty{1}{eV}=\qty{1,60e-19}{J}$ & Volume molaire (\qty{0}{\degreeCelsius} ; \qty{1,013}{bar}) & \qty{22,4}{L.mol^{-1}} \\
Kilowattheure & $\qty{1}{kW.h}=\qty{3,6e6}{J}$ & Volume molaire (\qty{20}{\degreeCelsius} ; \qty{1,013}{bar}) & \qty{24,0}{L.mol^{-1}} \\
Pression & $\qty{1}{bar}=\qty{1e5}{Pa}$ & Eau : masse volumique & \qty{1,00e3}{kg.m^{-3}} \\
Température & $T(\unit{K})=\theta(\unit{\degreeCelsius})+\num{273,15}$ & Eau de mer : masse volumique & \qty{1,03e3}{kg.m^{-3}} \\
Énergie pour chauffer l'eau & \qty{4,18}{kJ.kg^{-1}.\degreeCelsius^{-1}} & Célérité du son dans l'air & \qty{340}{m.s^{-1}} environ \\
Terre : masse, rayon & \qty{5,97e24}{kg} ; \qty{6,37e6}{m} & Lune : masse ; distance à la Terre & \qty{7,35e22}{kg} ; \qty{3,84e8}{m} \\
Soleil : masse & \qty{1,99e30}{kg} & Distance Terre-Soleil & \qty{1,50e11}{m} \\\bottomrule
\end{tabularx}}
\editionEtape{Quelques couples oxydant-réducteur}
{\small\begin{tabularx}{\linewidth}{@{}YYYY@{}}\toprule
$\ce{H+}/\ce{H2}$ & $\ce{Cu^{2+}}/\ce{Cu}$ & $\ce{Fe^{2+}}/\ce{Fe}$ & $\ce{Fe^{3+}}/\ce{Fe^{2+}}$ \\
$\ce{Zn^{2+}}/\ce{Zn}$ & $\ce{Ag+}/\ce{Ag}$ & $\ce{Al^{3+}}/\ce{Al}$ & $\ce{I2}/\ce{I-}$ \\
$\ce{MnO4-}/\ce{Mn^{2+}}$ & $\ce{O2}/\ce{H2O2}$ & $\ce{H2O2}/\ce{H2O}$ & $\ce{S4O6^{2-}}/\ce{S2O3^{2-}}$ \\\bottomrule
\end{tabularx}}
\editionEtape{Énergies de liaison moyennes, en phase gazeuse (\unit{kJ.mol^{-1}})}
{\small\begin{tabularx}{\linewidth}{@{}YYYYYYY@{}}\toprule
C--H & C--C & C--O & C$=$O (dans $\ce{CO2}$) & O$=$O & O--H & H--H \\\midrule
415 & 345 & 358 & 804 & 498 & 463 & 436 \\\bottomrule
\end{tabularx}}
\end{blocEdition}

\editionPage{donnees-utiles-2}{Tableau périodique}{Masses molaires et électronégativités}{Éléments des cinq premières périodes}
\begin{blocEdition}{donnees-utiles-2}
\fontsize{10}{13.5}\selectfont
\label{figedition-donnees-2}\schema{annexe-tableau-periodique}
\editionEtape{Utiliser l'électronégativité}
L'électronégativité $\chi$ traduit la capacité d'un atome à attirer vers lui les électrons d'une liaison covalente. Dans une
liaison entre deux atomes d'électronégativités différentes, l'atome le plus électronégatif porte une charge partielle
négative $\delta^-$, l'autre une charge partielle positive $\delta^+$ : la liaison est polarisée. On considère en général qu'une
liaison est polarisée lorsque l'écart d'électronégativité dépasse 0,4 environ ; c'est le cas des liaisons C--O, O--H, N--H ou
C--Cl, mais pas de la liaison C--H. Une molécule est polaire si les centres des charges partielles positives et négatives ne
coïncident pas : sa géométrie compte autant que ses liaisons.
\editionEtape{Spectroscopie infrarouge : nombres d'onde caractéristiques}
{\small\begin{tabularx}{\linewidth}{@{}Yl|Yl@{}}\toprule
O--H d'un alcool (liaison H) & 3200 -- \qty{3400}{cm^{-1}}, large & C$=$O d'un ester & 1735 -- \qty{1750}{cm^{-1}} \\
O--H d'un acide carboxylique & 2500 -- \qty{3200}{cm^{-1}}, très large & C$=$O d'un acide & 1700 -- \qty{1725}{cm^{-1}} \\
N--H d'une amine ou d'un amide & 3100 -- \qty{3500}{cm^{-1}} & C$=$O d'un aldéhyde, d'une cétone & 1650 -- \qty{1730}{cm^{-1}} \\
C--H & 2800 -- \qty{3100}{cm^{-1}} & C$=$O d'un amide & 1650 -- \qty{1700}{cm^{-1}} \\\bottomrule
\end{tabularx}}
\end{blocEdition}

\editionPage{couverture-programme}{Repères du programme}{Notions et capacités exigibles}{Programme de spécialité de première (BO spécial n°~1 du 22 janvier 2019) · Renvois aux chapitres}
\begin{blocEdition}{couverture-programme}
\fontsize{10}{13.5}\selectfont
Le programme est organisé en quatre thèmes, complétés par une partie transversale, « Mesure et incertitudes ». Les
tableaux suivent son plan et indiquent les chapitres du livre où chaque notion est traitée. Les capacités exigibles
détaillées sont reprises en tête de chaque chapitre, dans la rubrique « Objectifs ».
\editionEtape{Constitution et transformations de la matière}
{\small\begin{tabularx}{\linewidth}{@{}p{50mm}Yr@{}}\toprule
\textbf{Partie du programme} & \textbf{Notions principales} & \textbf{Ch.} \\\midrule
1.A Composition du système initial & masse molaire, quantité de matière, volume molaire, concentration ; absorbance, loi de
Beer-Lambert, dosage par étalonnage & 1, 2 \\
1.B Évolution d'un système chimique & oxydant, réducteur, demi-équations, réaction d'oxydoréduction ; avancement, tableau
d'avancement, réactif limitant & 3, 4 \\
1.C Quantité de matière par titrage & titrage direct, équivalence, repérage colorimétrique & 5 \\
2.A Structure et polarité des entités & schéma de Lewis, géométrie, électronégativité, polarité & 6 \\
2.B Cohésion, solubilité, miscibilité & interactions entre entités, dissolution, extraction liquide-liquide, savons & 7 \\
3.A Structure des entités organiques & formules, groupes caractéristiques, familles, nomenclature, spectroscopie infrarouge & 8 \\
3.B Synthèses organiques & étapes d'une synthèse, rendement, identification & 9 \\
3.C Énergie de la matière organique & combustions, pouvoir calorifique, énergie de liaison, enjeux environnementaux & 10 \\\bottomrule
\end{tabularx}}
\editionEtape{Mouvement et interactions}
{\small\begin{tabularx}{\linewidth}{@{}p{50mm}Yr@{}}\toprule
1. Interactions et champs & loi de Coulomb, loi de la gravitation universelle, champs électrique et de gravitation, lignes de champ & 11 \\
2. Fluide au repos & masse volumique, pression, force pressante, loi de Mariotte, loi fondamentale de la statique des fluides & 12 \\
3. Mouvement d'un système & vecteur variation de vitesse, lien entre $\Delta\vec v$ et la somme des forces & 13 \\\bottomrule
\end{tabularx}}
\end{blocEdition}

\editionPage{couverture-programme-2}{Repères du programme (suite)}{Notions et capacités exigibles}{Programme de spécialité de première · Renvois aux chapitres}
\begin{blocEdition}{couverture-programme-2}
\fontsize{10}{13.5}\selectfont
\editionEtape{L'énergie : conversions et transferts}
{\small\begin{tabularx}{\linewidth}{@{}p{50mm}Yr@{}}\toprule
1. Aspects énergétiques des phénomènes électriques & porteurs de charge, intensité, source de tension réelle, puissance, énergie,
effet Joule, rendement d'un convertisseur & 14 \\
2. Aspects énergétiques des phénomènes mécaniques & travail d'une force, théorème de l'énergie cinétique & 15 \\
& énergie potentielle de pesanteur, énergie mécanique, conservation et non-conservation & 16 \\\bottomrule
\end{tabularx}}
\editionEtape{Ondes et signaux}
{\small\begin{tabularx}{\linewidth}{@{}p{50mm}Yr@{}}\toprule
1. Ondes mécaniques & onde progressive, célérité, retard, onde périodique, période, longueur d'onde & 17 \\
2.A Images et couleurs & lentille mince convergente, relation de conjugaison, grandissement, œil & 18 \\
& synthèses additive et soustractive, couleur des objets, vision des couleurs & 19 \\
2.B Modèles ondulatoire et particulaire & domaines des ondes électromagnétiques, photon, niveaux d'énergie, spectres & 20 \\\bottomrule
\end{tabularx}}
\editionEtape{Mesure et incertitudes}
{\small\begin{tabularx}{\linewidth}{@{}p{50mm}Yr@{}}\toprule
Partie transversale & variabilité de la mesure, histogramme, moyenne, écart-type, incertitudes-types de type A et B, écriture
du résultat, comparaison qualitative à une valeur de référence & 21 \\\bottomrule
\end{tabularx}}
\editionEtape{Capacités numériques du programme}
{\small\begin{tabularx}{\linewidth}{@{}Yr@{}}\toprule
Représenter l'histogramme associé à une série de mesures à l'aide d'un tableur & 21 \\
Déterminer la composition de l'état final d'un système siège d'une transformation totale & 4 \\
Étudier la relation approchée entre la variation du vecteur vitesse et la somme des forces & 13 \\
Effectuer le bilan énergétique d'un système en mouvement & 16 \\
Représenter un signal périodique ; simuler la propagation d'une onde périodique & 17 \\\bottomrule
\end{tabularx}}\par\smallskip
Le livre propose d'autres programmes courts dans chaque chapitre : chacun est repéré par l'étiquette « Capacité numérique ».
\end{blocEdition}
